
\emergencystretch=2em
\raggedbottom
\hbadness=5000
\vbadness=5000

\maketitle


\begin{abstract}
Can the physical selection of observable algebras account for gravitational
dynamics, rather than merely change what an observer can reconstruct?
We formulate this question as an adiabatic evolution of selected localization
structures and a common-action consistency requirement. Modular stability
alone does not select the algebra. A supplied half-sided modular family
and holographic dictionary give a conditional linearized gravity route.
In the same-spacetime conformal scalar model of Consciousness~II, the
improved stress and its Wald surface coefficient connect accessible
preparations to an all-ball test of the Einstein response. The test still
requires modular consistency; changing a stress-containing algebra at
fixed matter history does not change its first-order sourced Einstein tensor.

The source side is constructive. Calibrated correction, selected projections
and learned protection change stress through actual interactions. In a
specified Jackiw--Teitelboim theory, an acquired quantum record drives a
material force and radiation, whose response changes later operations.
The finite mean-action return retains old references, clock energy,
support work and recoil. Separate spatial and full-mode bulk constructions
extend the source account without identifying their distinct models.

Maintenance supplies a possible selection pressure. For complete entropy
removal at rate $\dot s$ through a thermal chiral exhaust of fixed capacity
$c_{\rm ex}$, the heat rate obeys
$\dot Q\ge k_BT\dot s+3\hbar\dot s^2/(\pi c_{\rm ex})$.
Pointer alignment lowers the complete renewal debit, and matched horizon
transport reads its price. Sustained renewal and joint sharp exhaust with
persistent causal return remain open. The proposed gravitational mechanism
must connect formed access, a conserved total source and geometric
consistency in one dynamics; neither a visibility filter nor a finite
response in an assigned gravity theory supplies that connection.
\end{abstract}

\section{The problem: what makes selection gravitational?}
\label{sec:intro}

An apparatus changes which distinctions it can preserve and which operations
it can execute. This can reorganize a reconstructed geometry. It can also
change matter through the selecting interaction. A theory of gravity based
on observable-algebra selection must connect these two effects: the
localization structure selected by physical histories must constrain a
geometric law, and the same histories must supply the source and respond to it.

The central proposal is that gravitational dynamics describes a slow
evolution of the selected geometric algebra, not just motion of matter
within a fixed localization structure. This is a mechanism conjecture,
developed in Section~\ref{sec:technical}, rather than an identification of
every algebra change with curvature. Its strongest obstruction is already
visible in a conformal model: two selected algebras that retain the same
stress readings can give different relative entropies while fixing the
same matter-sourced Einstein tensor. A physical change of localization
must therefore be distinguished from a change of description.

HAFF-A, B and E relate channels, stable content and actual records
\cite{Liu2026PaperA,Liu2026PaperB}. Foundation~I constructs a stable hull
relative to a readable seed \cite{FoundationI}. These results make parts
of algebra selection precise; they do not determine a spacetime metric
or a universal gravitational dynamics. Here we ask what further relation
would make that selection gravitational.

Sections~\ref{sec:factorization}--\ref{sec:modular} state the algebraic
question, its central conjecture and the modular constraints. The
holographic route needs an explicit region family and dictionary.
Section~\ref{sec:common-geometry} gives a different conditional bridge
in the same spacetime: one conformal scalar source, its surface charge,
and its all-ball geometric test. That bridge identifies which law is
compatible with the supplied consistency condition, while leaving the
physical origin of that condition to the common-action conjecture.

Sections~\ref{sec:source} and \ref{sec:physical-return} then test the
source and return edges through actual interactions in a specified
two-dimensional JT theory. The decisive finite mechanism is an acquired
record that changes a material source, not a command label attached to a
separate emitter. Section~\ref{sec:gap} prices the full maintained state.
Section~\ref{sec:sustainable} asks whether the source, renewal and
informative return can persist together. Long controller, clock,
support and source proofs are in the appendices.

Table~\ref{tab:scope-ledger} separates these constructions. Its rows are
not pieces of one already realized apparatus. In particular, a supplied
JT response cannot establish the higher-dimensional consistency law,
and a modular constraint cannot manufacture a controller.

\begin{table}[p]
\centering\small
\begin{tabular}{P{3.0cm}P{5.3cm}P{5.6cm}}
\toprule
Object and supplied inputs & Result & Remaining interface\\
\midrule
Modular family, dictionary and ball laws &
Conditional translations and linearized geometric consistency
(Sections~\ref{sec:factorization}--\ref{sec:modular}) &
Formation of the candidate family and its geometric dictionary.\\[.6ex]
Conformal scalar, source reference, collar prescription and MC &
One prepared source, matched Wald coefficient and all-ball Einstein test
(Section~\ref{sec:common-geometry}) &
Physical origin of MC; apparatus source and finite-resource readings.\\[.6ex]
Calibrated field operations and a JT action &
Derived outgoing source, retarded response and finite echo/contact
readout (Section~\ref{sec:source}; Appendix~\ref{app:boundary-return}) &
Total apparatus stress and physical references; the action is an input.\\[.6ex]
Quantum core with classical mechanics and mean-action source &
Finite retained-history return and strong-domain perturbation bounds
(Section~\ref{sec:physical-return}; Foundation~III) &
Fixed-window mean-action result; protected content, old references and
initial work remain charged. No quantum stress tensor for the scalar
fields or support is supplied.\\[.6ex]
Spatial quadratic fields and real-quartet bulk &
Work-corrected full-mode stress matching, a finite spatial return,
and a separately constrained quartet interior &
Joining their receivers, moving endpoints, preparation and thermal
ports in one material model.\\[.6ex]
Complete retained entropy and paid finite carrier &
Long-window heat floor and neutral-sector sorting of a resettable
source (Section~\ref{sec:gap}; Proposition~\ref{prop:pointer-sort}) &
Actual adaptive source concentration, protected-content compatibility
and local delayed readout.\\[.6ex]
One reused body over growing runs &
Conjecture~\ref{conj:gravitating} jointly requires complete-output
sharpness, normalized window contrast and causal transition gain &
Sustained inlet, support, sensing and resource account in that same
dynamics; none is supplied by sorting alone.\\
\bottomrule
\end{tabular}
\caption{Objects, results and interfaces. Finite constructions retain their
stated inputs; the joint growing-window target is a conjecture.}
\label{tab:scope-ledger}
\end{table}

\section{The Factorization Problem}
\label{sec:factorization}

\subsection{No canonical tensor factorization}

A Hilbert space of dimension $d=d_1d_2$ admits a tensor factorization,
but its Hilbert structure alone does not select one. Commuting observable
algebras can induce the subsystem structure~\cite{Zanardi2001,Zanardi2004}.
Which of those algebras an apparatus can use is a physical question about
its couplings and resources.

\subsection{Coarse-graining and effective descriptions}

Paper~A~\cite{Liu2026PaperA} gives the following conditional separation.

\begin{proposition}[Conditional coarse-graining separation]
\label{prop:inequiv}
Let $\Phi_1\ne\Phi_2$ be finite-dimensional coarse-graining channels with a common output space. Their outputs differ for generic pure inputs. Entropy profiles differ only under the additional nondegeneracy hypothesis stated in Paper~A, and metric inequivalence follows only when a supplied reconstruction dictionary separates those profiles modulo diffeomorphisms.
\end{proposition}

Channel inequivalence therefore establishes a difference of effective descriptions, not by itself a difference of spacetime metrics. Paper~A also decides, for qubit reductions, exactly when two channels give the same output entropy on every pure input, and it constructs one global state, one family of local noise channels and one fixed reconstruction rule that yield non-isometric finite metrics. Metric inequivalence of spacetime geometries remains a conditional model claim throughout this paper.

\subsection{Geometry from entanglement}

In a holographic model the Ryu--Takayanagi formula~\cite{RyuTakayanagi2006} relates an extremal area to boundary entropy,
\begin{equation}
S_A = \frac{\text{Area}(\gamma_A)}{4 G_N \hbar},
\end{equation}
and motivates an entanglement account of connectivity~\cite{VanRaamsdonk2010}.
Such a reconstruction requires its dictionary. Outside a specified model,
the accessible data determine operator relations rather than a spacetime metric.

\section{A Structural Distinction: Gauge Forces and Gravity}
\label{sec:distinction}

Gauge interactions evolve matter and connection fields within a supplied localization structure. The algebraic proposal asks whether gravity instead describes how physical correlations, accessible operations and their stability select that structure. With a geometric dictionary, changing those relations can change effective distances and causal structure. The source constructions below test a necessary interface: changed access reaches geometry through a specified matter process.

If the proposal holds, quantum gravity must describe selection of localization as well as fields within it. This is not an explanation of why quantization fails. Background independence~\cite{Rovelli2004,Smolin2006} asks for covariance under changes of selected geometric structure, not equality of the metrics assigned to operationally different algebras. Selection is constrained by interactions and stability, rather than subjective choice~\cite{Liu2026PaperB}; changed knowledge affects stress only through an operation.

A common reconstruction dictionary gives every matter representation in one selected geometric sector the same metric. This conditional universality is not a derivation of free fall or equality of inertial and gravitational mass. Degrees outside the sector cannot be declared geometrically irrelevant while changing its metric: a dark-sector application would need an enlarged source algebra and explicit backreaction map. No dark-matter or dark-energy conclusion follows here.

\section{Algebra selection as a gravitational mechanism}
\label{sec:technical}

\subsection{Setup}

We separate the operator-algebraic and finite-dimensional regimes result by result. Modular statements use a Type~III$_1$ ambient factor with a faithful normal state and explicit cyclic and separating hypotheses. Density matrices, traces, the quantum Fisher information metric and the qubit chain are finite-dimensional statements. The Rindler example is an algebraic quantum field theory check. Theorem~\ref{thm:linearized} uses the holographic regime; Section~\ref{sec:common-geometry} uses a distinct same-spacetime conformal model. The JT model of Section~\ref{sec:source} is a supplied two-dimensional gravity theory.

\begin{definition}[Selected geometric algebra]
\label{def:accessible}
A \textbf{selected geometric algebra} $\mathcal{A}_{\mathrm{geo},\mathbf{c}}$ is a von Neumann subalgebra chosen as input to a specified algebra-to-geometry reconstruction model. It is not Paper~E's probe-reachable effect system, pointer or record algebra, or environmental fixed set. A coarse-graining model must declare the channel or inclusion relating any of those objects to $\mathcal{A}_{\mathrm{geo},\mathbf{c}}$.
\end{definition}

No canonical tensor factorization is assumed. State restrictions and relative-entropy data on $\mathcal{A}_{\mathrm{geo},\mathbf{c}}$ can be associated, as in holography and tensor networks, with an effective distance structure. That structure depends on the choice of $\mathcal{A}_{\mathrm{geo},\mathbf{c}}$, on the stability of correlations under coarse-grained dynamics and on the redundancy of the encoding; it is an effective description valid in the context the algebra defines.

\subsection{Central conjecture}

\begin{conjecture}[Gravity as adiabatic algebra evolution]
\label{conj:main}
Gravitational dynamics and the \textbf{adiabatic flow} of selected geometric algebras $\mathcal{A}_{\mathrm{geo},\mathbf{c}}(t)$ arise from one physical selection process:
\begin{enumerate}
    \item the global state evolves unitarily, $|\Psi_U(t)\rangle = U(t)|\Psi_U(0)\rangle$;
    \item a model-dependent rule selects a geometric candidate $\mathcal{A}_{\mathrm{geo},\mathbf{c}}(t)\subset\mathcal M_U$ at each time;
    \item as the state evolves, the selected candidate may shift, $\mathcal{A}_{\mathrm{geo},\mathbf{c}}(t) \to \mathcal{A}_{\mathrm{geo},\mathbf{c}}(t + dt)$;
    \item the interaction that changes the selected algebra also determines the complete physical source and the compatible localization data. A change of algebra alone, with the matter history and source readings held fixed, supplies no additional first-order sourced Einstein tensor in the common-source regime below. A changed sourced response must be carried by the changed total source under that regime's field-law conditions.
\end{enumerate}
The proposed distinction from ordinary matter feedback is that the same dynamics selects the localization structure in which the source and response are represented, rather than merely evolving matter on a permanently supplied stage. Evolution within a fixed selected localization can describe gauge dynamics, but neither unitarity nor a moving algebra by itself identifies an interaction as gauge or gravitational. The selection claim must be established through the common-action conditions below; finite sourced return alone does not establish it.
\end{conjecture}

The proposed flow needs a physical selector. A readable history changes
which controls and correlations remain stable; those controls determine a
selected hull and the matter interactions that sustain it. The changing
hull must then admit compatible geometric region data, rather than merely
receive an arbitrary distance assignment. In an adiabatic sector, the
relaxation time of the selected structure must be short compared with its
change time. This separation is a condition to establish in a model, not
a consequence of writing $\mathcal A(t)$.

Two explanations compete here. In the reconstruction explanation, the
ambient geometry and complete source stay fixed, while smaller algebras
resolve less of them. In the dynamical explanation, history changes an
actual coupling or preparation, changes the complete source, and changes
later physical readings. The latter has a finite realization below; the
stronger claim that the localization law itself is selected requires
Section~\ref{sec:source-principle}. Algebra variation alone distinguishes
neither explanation.

A useful test compares histories under one action and one source
calibration. Erasing command access before execution should remove the
record-mediated part of the source; erasing a log after execution should
not. If two selected structures contain the same total stress and surface
data, a first-order difference of their sourced Einstein tensors would
contradict the conformal bridge in Section~\ref{sec:common-geometry}.
A difference only in relative entropy supports a reconstruction claim
unless a corresponding physical intervention is demonstrated. Conversely,
a history-dependent source that fails the same all-ball response test
does not realize the proposed gravitational flow, even if it produces a
detectable JT signal. Rapid algebra transitions lie outside the adiabatic
claim; whether they admit a different geometric regime is open.

\subsection{Comparison with quantum information geometry}

For comparison consider the quantum Fisher information metric (QFIM) on a smooth finite-dimensional family of faithful effective states $\rho(\lambda)$~\cite{Petz1996,Bengtsson2006},
\begin{equation}
g_{\mu\nu}(\lambda) = \frac{1}{2} \text{Tr}\left( \rho(\lambda) \{ L_\mu, L_\nu \} \right),
\qquad
\partial_\mu \rho = \frac{1}{2}(\rho L_\mu + L_\mu \rho).
\label{eq:qfim-d}
\end{equation}
It is coordinate independent, reduces to the classical Fisher metric in commuting limits, and measures how hard nearby states are to distinguish. It is also positive semidefinite, whereas a Lorentzian spacetime metric is indefinite and carries causal structure. Nothing in Eq.~\eqref{eq:qfim-d} supplies the negative direction, so the QFIM is a metric on model space and an analogy, not the spacetime metric. A change $\mathcal{A}_{\mathrm{geo},\mathbf{c}} \to \mathcal{A}_{\mathrm{geo},\mathbf{c}} + \delta\mathcal{A}_{\mathrm{geo},\mathbf{c}}$ may change this state-space metric; calling the change a spacetime perturbation requires an additional Lorentzian reconstruction map. In AdS/CFT, relative entropy and the entanglement first law constrain bulk perturbations~\cite{Lashkari2014,Faulkner2014}; those results use the holographic dictionary and do not make an arbitrary boundary QFIM the bulk metric.


\section{Algebraic Foundation: Modular-Stable Candidates}
\label{sec:modular}

This section places one candidate class in Tomita--Takesaki modular theory. It does not translate measurement accessibility or environmental redundancy into modular conditions, and it does not select a unique algebra.

Standard modular definitions and their two checks are in Appendix~\ref{app:modular-details}.

\subsection{Modular definition of candidate algebras}
\label{sec:modular-def}

Let $\mathcal M_U$ be a Type~III$_1$ factor in standard form on $\mathcal H_U$, let $|\Psi_U\rangle$ be cyclic and separating, and write $\omega(\cdot)=\langle\Psi_U|\cdot|\Psi_U\rangle$. Then $\omega$ is faithful and normal and its ambient modular group $\sigma_t^{\omega,U}$ exists. Let $\{\mathcal E_t\}_{t\ge0}$ be a normal unital completely positive (UCP) semigroup acting in the Heisenberg picture on $\mathcal M_U$.

\begin{definition}[Modular-stable geometric candidate]
\label{def:modular-accessible}
A von Neumann subalgebra $\mathcal{A}_{\mathrm{geo},\mathbf{c}}\subset\mathcal M_U$ is a \textbf{modular-stable geometric candidate} if:
\begin{enumerate}
\item[\textup{(U1)}] \textbf{Faithful restriction.} $\omega|_{\mathcal{A}_{\mathrm{geo},\mathbf{c}}}$ is faithful and normal. In the faithful ambient setting this holds for every unital subalgebra, so it is a regularity condition, not a selector.
\item[\textup{(U2)}] \textbf{Ambient modular stability.} $\sigma_t^{\omega,U}(\mathcal{A}_{\mathrm{geo},\mathbf{c}})=\mathcal{A}_{\mathrm{geo},\mathbf{c}}$ for all real $t$.
\item[\textup{(U3)}] \textbf{Semigroup stability.} $\mathcal E_t(\mathcal{A}_{\mathrm{geo},\mathbf{c}})\subseteq\mathcal{A}_{\mathrm{geo},\mathbf{c}}$ for all $t\ge0$.
\end{enumerate}
\end{definition}

\begin{remark}[Typed relation to the measurement objects]
\label{rem:translation}
Conditions (U1)--(U3) do not derive Paper~E's implementable effects $\mathrm{Eff}$, their operator system $\mathrm{Eff}_{\mathrm{span}}$ or the pointer algebra $\mathcal P_{\mathrm{env}}$, and setwise semigroup stability is not the exact fixed set $\operatorname{Fix}(\mathcal E)$. A downstream model may specify a channel or inclusion from those objects into $\mathcal{A}_{\mathrm{geo},\mathbf{c}}$; without it the relation is motivational. Environmental redundancy does not imply faithfulness of a state restriction. A causal restriction map and a Lindblad evolution are different types, and each model must supply its own $\mathcal E_t$ and verify invariance.
\end{remark}

\subsection{Selection is not settled by these conditions}
\label{sec:modular-uniqueness}

\begin{remark}[No uniqueness from (U1)--(U3)]
\label{thm:uniqueness}
Conditions (U1)--(U3), even with only scalars fixed by both dynamics, do not determine a unique candidate. There are four independent obstructions.
\begin{enumerate}
\item Uniqueness of the modular flow for a fixed pair $(\mathcal A,\omega)$ does not select the algebra $\mathcal A$.
\item A normal completely positive map may preserve two subalgebras separately without preserving their join, because it is not multiplicative. Join preservation holds for modular automorphisms but not for a general $\mathcal E_t$; it does hold for subalgebras of the fixed algebra of a state-preserving dynamics, on which $\mathcal E_t$ acts as the identity~\cite{FoundationI}.
\item Pointwise ergodicity constrains fixed elements, while (U2)--(U3) are setwise conditions; no implication to a unique maximal element of the lattice was proved.
\item With a faithful ambient state, $\mathcal M_U$ itself satisfies (U1)--(U3), so a bare maximality rule selects the ambient algebra unless an independent seed, properness, read-out or minimal-hull condition is supplied.
\end{enumerate}
Replacing $\mathcal B(\mathcal H_U)$ by a faithful Type~III factor repairs the existence of the ambient modular flow but none of these selection failures. The class of candidates may contain several members that are not unitarily equivalent, and classifying it requires additional structure. A nontrivial selection principle could start from a foreground seed and construct a minimal invariant hull, or impose an explicit read-out constraint. Foundation~I carries out the first route for a supplied seed: under state-preserving dynamics the hull of a seed in the fixed algebra lies between that seed and the fixed algebra. For a minimal compact symmetry, it equals the seed-stabilizer fixed algebra when that algebra is preserved by the hull generators and the hull contains the group-fixed algebra~\cite{FoundationI}. These hypotheses establish uniqueness relative to the seed; the seed's physical origin is a further question.
\end{remark}

\subsection{Connection to linearized gravity}
\label{sec:linearized}

\begin{theorem}[Entanglement first law in a density-matrix setting]
\label{thm:first-law}
Let $\rho_0$ be a faithful finite-dimensional density matrix and
$\rho_\varepsilon=\rho_0+\varepsilon\,\delta\rho+O(\varepsilon^2)$ a normalized
perturbation, with $K_0=-\log\rho_0$. Then $\delta S=\delta\langle K_0\rangle$.
An infinite-dimensional extension requires entropy differentiability and control
of the modular-energy variation; faithfulness alone is insufficient. For a Type~III
local algebra the statement must instead use Araki relative entropy.
\end{theorem}

\begin{proof}[Proof sketch]
$S(\rho_\varepsilon\|\rho_0)=-S(\rho_\varepsilon)+\langle K_0\rangle_{\rho_\varepsilon}$ plus a constant is nonnegative and stationary at $\varepsilon=0$.
\end{proof}

Conjecture~\ref{conj:main} posits a family $\mathcal{A}_{\mathrm{geo},\mathbf{c}}(\lambda)$. Modular theory treats two kinds of perturbation differently. For a state perturbation on a fixed algebra the Connes cocycle $u_t=(D\omega':D\omega)_t$ relates the automorphism groups by
$\sigma_t^{\omega'}(A)=u_t\sigma_t^\omega(A)u_t^*$~\cite{Connes1994,Takesaki2003}. This is not a left-multiplication formula for the two full modular operators. If the algebra family is generated by the positive translation of an HSMI,
$\mathcal A_{\mathrm{geo},\mathbf c}(a)=e^{iaP}\mathcal A_{\mathrm{geo},\mathbf c}(0)e^{-iaP}$,
and the vacuum is fixed by this translation, its full modular generator obeys
$K(a)=e^{iaP}K(0)e^{-iaP}$. On a common smooth domain,
\begin{equation}
\label{eq:deltaK-translation}
\frac{dK(a)}{da}\Big|_{a=0}=i[P,K(0)] .
\end{equation}
Equation~\eqref{eq:deltaK-translation} concerns an HSMI-generated family and a derivative
of matrix elements on that domain, not arbitrary subalgebra changes or an operator-norm
Taylor expansion. The one-sided CHM charge used for a ball below is distinct from the
full modular generator; their relevant local-unitary expectation differences agree,
but the operators themselves do not.

\begin{theorem}[Linearized gravity from entanglement, {\cite{Faulkner2014}}]
\label{thm:linearized-Faulkner}
In a holographic CFT dual to Einstein gravity in asymptotically AdS spacetime, the entanglement first law $\delta S = \delta\langle K_{\mathrm{mod}}\rangle$ for all boundary balls is equivalent to the linearized bulk Einstein equations $G_{\mu\nu}^{(1)} + \Lambda g_{\mu\nu}^{(1)} = 8\pi G_N\, T_{\mu\nu}^{(1)}$.
\end{theorem}

The derivation uses the ball modular Hamiltonian~\cite{CasiniHuertaMyers2011}
\begin{equation}
K_{\mathrm{mod}}^B = 2\pi \int_B d^{d-1}x\, \frac{R^2 - |\mathbf{x} - \mathbf{x}_0|^2}{2R}\, T_{00}(x),
\end{equation}
together with the Ryu--Takayanagi formula and the relation between bulk and boundary modular Hamiltonians~\cite{JLMS2016}; complementary arguments are due to Jacobson~\cite{Jacobson2016}.

\begin{theorem}[Conditional holographic matching for an HSMI-generated family]
\label{thm:linearized}
Let $\mathcal{A}_{\mathrm{geo},\mathbf{c}}$ be a modular-stable candidate in a holographic CFT vacuum representation, and assume:
\begin{enumerate}
\item[(H1)] the AdS/CFT correspondence holds;
\item[(H2)] the Ryu--Takayanagi formula and its quantum corrections hold;
\item[(H3)] the JLMS relation between bulk and boundary modular Hamiltonians holds;
\item[(H4)] the compared algebra family is an explicitly supplied HSMI/ball-region family whose generator $P$ implements the boundary-region variation used in the holographic first-law argument;
\item[(H5)] the regulated entropy variation, or its relative-entropy replacement, satisfies the hypotheses of the Faulkner et al.\ theorem for all required balls.
\end{enumerate}
Then Eq.~\eqref{eq:deltaK-translation} describes the modular-generator variation within that family, and the holographic first law matches the linearized Einstein equations by the imported result:
\begin{align}
 \text{HSMI/ball family (H4)}&\Longrightarrow\delta K_{\mathrm{mod}}
 \xrightarrow{\text{(H5)}}\delta S=\delta\langle K_{\mathrm{mod}}\rangle,
 \nonumber\\
 \text{(H1)--(H3)}&\Longrightarrow
 \delta(G_{\mu\nu}+\Lambda g_{\mu\nu})=8\pi G_NT_{\mu\nu}^{(1)}.
\end{align}
The theorem does not show that a generic stability-selected algebra variation is an HSMI translation, nor that HAFF algebra flow is gravitational dynamics.
\end{theorem}

\subsection{The circularity issue}

Jacobson's thermodynamic derivation~\cite{Jacobson1995} presupposes a background geometry, while HAFF wants geometry to emerge from algebra. Three strategies are discussed in the literature. (A) In the algebraic programme, wedges are defined through modular inclusions without a background metric, and Theorem~\ref{thm:Wiesbrock} recovers a positive translation from an HSMI; constructing a sufficient net and showing that the candidates of this paper form such inclusions remain open tasks. (B) A bootstrap starts from a seed geometry, derives linearized equations from the first law, updates the geometry and iterates. (C) In holography the boundary theory supplies the algebra and the bulk geometry is derived from boundary data. Theorem~\ref{thm:linearized} uses strategy (C) and assumes the HSMI/ball interface. A geometry-free reconstruction from the present candidate class remains open.


\section{One source and a same-spacetime geometric test}
\label{sec:common-geometry}

The holographic matching above is not the only conditional route.
Consciousness~II \cite{PaperCII} uses a free massless conformally coupled
complex scalar in $D=d+1\ge3$ dimensions. Its ball algebras, physical
translation generator and vacuum are supplied in Minkowski spacetime.
This matters: they are not the finite pointer algebra of a JT reader.
What can be combined in this model is the prepared source, its modular
charge and a geometric consistency test in that same spacetime.

In units $\hbar=c=1$, the scalar curvature coupling is
$-\xi\mathcal R\Phi^2/2$, with
$\Phi^2=\phi_1^2+\phi_2^2$ and $\xi=(d-1)/(4d)$.
Metric variation gives the conformally improved stress.
For a ball $B=B(x_0,R,u)$, write
\begin{align}
 K_B&=2\pi\int_B w_R T_{uu}\,d^dx,\qquad
 w_R=\frac{R^2-|x-x_0|^2}{2R},\nonumber\\
 H_B&=K_B+Q_B,\qquad
 Q_B=2\pi\xi\int_{\partial B}:\!\Phi^2\!:\,dA .
 \label{eq:common-ball-data}
\end{align}
These are vacuum-subtracted expectation data with a fixed scalar boundary
prescription. Surface readings mean controlled collar limits.
$H_B$ is the physical one-sided energy, $K_B$ its bulk charge;
neither is the full modular generator $-\log\Delta_B$.
Keeping the same curvature coupling with Einstein coefficient
$1/(16\pi G)$ gives Wald density $1/(4G)-2\pi\xi\Phi^2$.
The surface coefficient is therefore fixed by the same matter action,
not chosen separately for every reader.

\begin{proposition}[A common-source conditional bridge]
\label{prop:common-source-bridge}
Use the conformal scalar representation and surface prescription above.
Let $\delta T$ and $h$ be smooth source and metric perturbations on an
open domain. For the intrinsic geodesic ball of fixed proper radius set
\[
 \delta F_B[h]=\delta\left(A(\partial B)-\frac{d-1}{R}V(B)\right),
 \qquad \eta_A>0 .
\]
Then the following statements hold.
\begin{enumerate}
\item The bulk condition
\begin{equation}
 \eta_A\delta F_B[h]+\delta\langle H_B\rangle
                 -\delta\langle Q_B\rangle=0
 \label{eq:D-MC}
\end{equation}
for all centres, radii and rest frames is equivalent to
$G^{(1)}[h]=(2\pi/\eta_A)\delta T$.
It suffices to test the leading small-ball term at every centre and
an open set of rest frames.
\item At $\eta_A=1/(4G)$ the same condition is
$\mathcal B_B+\delta\langle H_B\rangle=0$, where
\[
 \mathcal B_B=\delta S_{\rm W}[\partial B]
            -\eta_A\frac{d-1}{R}\delta V(B)
            =\eta_A\delta F_B-\delta\langle Q_B\rangle
\]
at linear order about the vacuum reference. The source normalization is
$8\pi G$ in both forms.
\item If selected ball subalgebras have vacuum-preserving normal
conditional expectations and contain the bounded stress readings and
the required surface readings, restricting one fixed physical state
to them leaves Eq.~\eqref{eq:D-MC} and its sourced Einstein tensor
unchanged. A metric additionally requires common gravitational initial
or boundary data and gauge identification.
\item Suppose a preparation has a smooth source with integrable energy
density of strictly positive total energy, and that the source and
response Fourier transforms belong to the tempered tensor-valued Radon
measure class of Consciousness~II. If consistency is tested at all
centres and an open frame set, its nonzero-frequency data identify the
canonical intrinsic response up to one normalization at each tested
radius, within that paper's intrinsic, gauge-invariant, isotropic,
compactly supported linear response class.
\end{enumerate}
\end{proposition}
\begin{proof}
The exact flat ball identity is
$\delta F_B=-\int_Bw_RG^{(1)}_{uu}$ \cite[Theorem on the exact
identity]{PaperCII}. Subtracting $Q_B$ from $H_B$ in
Eq.~\eqref{eq:common-ball-data} gives the other integrand
$2\pi\delta T_{uu}$. The leading coefficient of their sum is a nonzero
multiple of $R^{d+1}(-\eta_AG^{(1)}+2\pi\delta T)_{uu}$.
An open frame set recovers the symmetric tensor; the field equation
conversely makes every ball residual zero.
The curvature derivative of
$(1/(16\pi G)-\xi\Phi^2/2)\mathcal R$ gives the displayed Wald density,
so its scalar part is exactly $-Q_B$ at this order.
Conditional expectations preserve the vacuum modular restriction and
the common stress readings, proving the third statement.
For the fourth, a Radon Fourier measure supported only at zero would
give a constant source. Integrability forces its energy density to
vanish, contrary to the positive energy. The nonzero-frequency
identification theorem of Consciousness~II now applies with its stated
response hypotheses. These are consequences of that model's results,
not an independent derivation of modular consistency.
\end{proof}

The preparation hypotheses are nonempty. In the same scalar theory,
a local Weyl pulse gives a smooth improved stress with positive finite
energy and a rapidly decreasing mass-shell convolution measure.
A neutral pulse
$U_\theta=\exp[i\theta(\phi_1(f)^2+\phi_2(f)^2)/2]$
belongs to a proper symmetry-fixed stable hull when its bounded seed is
permitted. It too has a smooth positive-energy nonzero-frequency source
\cite[Sections on prepared sources and selected algebras]{PaperCII}.
The permission to execute that pulse is an input; containing stress
observables alone need not contain its generator.

There is a direct first-order comparison using this same preparation.
Fix $\theta\ne0$ and take the one-sided family
$\psi_\epsilon=(1-\epsilon)\omega+\epsilon\omega_{U_\theta}$,
$0\le\epsilon\le1$. Its source and collar tangents are exactly the
neutral pulse's finite-amplitude differences. Choosing the
constraint-compatible linear metric response to that conserved
source makes Eq.~\eqref{eq:D-MC} hold for every ball, with the same
Wald coefficient. The mixture can be prepared by a recorded random
choice of the permitted pulse or identity; that choice and the
apparatus are charged. The prepared matter source is not the whole
apparatus source when controls occupy the tested domain.
This common-model comparison uses the Einstein coupling and
gravitational constraint data as inputs, not outputs of the pulse.

A bounded non-scalar accessible observable with a permitted finite
translation interval can instead be time-smoothed to give a bounded
pulse of strictly positive finite energy. Local stress regularity and
the Radon condition still require separate proof in that general route.
Exact positive-energy translation permission on a nontrivial separated
finite-dimensional algebra is impossible in the unique-vacuum setting;
a finite clock is not that translation generator.

The fixed-source statement also survives finite adaptive reading
within closed permissions. If $E:\mathcal A\to\mathcal B$ is a normal
conditional expectation and every history-conditioned Heisenberg
instrument preserves $\mathcal B$, initial states $\rho$ and
$\rho\circ E$ give the same record probabilities and, on positive-probability
histories, conditional bounded source readings in $\mathcal B$ \cite{PaperCII}.
For each history, the pulled-back record effect and source observable
stay in $\mathcal B$, where $E$ is identity. Unbounded readings need
absolute expectation convergence of spectral truncations.
This permits real feedback that changes a common source; it forbids
that same restricted feedback from recovering a discarded initial
distinction. Enlarged controls or side information change the premise.

The bridge constrains the law rather than merely measuring one response:
within its intrinsic response class, a nonzero spectral witness removes
the higher-spatial-derivative ambiguity left by a homogeneous source.
It does not select an entire metric from matter data; homogeneous
gravitational solutions and common boundary data remain.
Raw scalar $H_B$ cannot replace the subtracted charge on all small
balls: its surface term is $O(R^{d-1})$, while smooth bulk stress and
geometry begin at $O(R^{d+1})$.
Nor is the Wald coefficient match a modular first law for an arbitrary
curved diamond. The area coefficient, region family, source reference,
preparation permissions and Eq.~\eqref{eq:D-MC} are still supplied.

\subsection{Conservation before geometric interpretation}
\label{sec:source-reaction}

A source eligibility rule cannot remove unread stress at will.
Q-RAIF~NV \cite[Section on physical selection and reaction]{QRNV}
proves a limited obstruction: on a connected flat domain of $D\ge2$,
a filter $T\mapsto a(x)T+b(x)\operatorname{tr}(T)g$ that conserves
every smooth symmetric conserved $T$ has constant $a$ and $b=0$.
If it also kills every constant vacuum source, it is zero.
This is a statement about that local algebraic filter family, not a
no-go theorem for additional-field or nonlocal gravity models.

The same paper exhibits the appropriate alternative at the classical
level. For an actual coupling field $\chi$ and canonical matter
$\mathcal L_m=-\frac12(\nabla\psi)^2-U(\psi)$, let
\[
 S=\int\sqrt{-g}\left[-\tfrac12(\nabla\chi)^2-V(\chi)
                                      +f(\chi)\mathcal L_m\right],
 \qquad f>0 .
\]
On both field equations,
\[
 \nabla^\mu(fT^m_{\mu\nu})=\mathcal L_m\partial_\nu f
                      =-\nabla^\mu T^\chi_{\mu\nu}.
\]
The changed material coupling carries a compensating physical stress.
This standard action-and-reaction identity is not a quantum total-source
construction or a history-derived selector. It also does not cancel a
vacuum term: shifting $\mathcal L_m$ by $-\rho$ changes both the total
stress and the $\chi$ equation.

There is a concrete incompatibility to resolve before combining this
mechanism with Proposition~\ref{prop:common-source-bridge}.
A varying $f(\chi)$ changes the matter theory, so its ball modular
Hamiltonian is not automatically the free conformal one.
One may compare the prepared source in a causal domain where that
interaction has ended and the specified conformal dynamics applies,
with selector stress absent there or explicitly accounted for; or one
must derive the interacting modular response anew.
A separately prepared switching field plus the free CHM formula is
not a common-model proof.

\subsection{The second-order test}

At the next order the source must follow the changed geometry.
Consciousness~II writes
$g_\epsilon=\eta+\epsilon h+\epsilon^2j+\cdots$ and
$T_\epsilon=\epsilon t_1+\epsilon^2t_2+\cdots$.
Its corrected small-ball test yields
\[
 G^{(1)}[j]+G^{(2)}[h,h]=(2\pi/\eta_A)t_2
\]
only after adding the slice term
$-\frac{V_d}{2(d+2)}R^{d+1}
[(\operatorname{tr}K)^2-K_{ij}K^{ij}]$ to the area--volume functional,
or choosing a geodesic slice with $K(p)=0$.
The source coefficient obeys the expanded covariant Ward identity,
including its metric variation and changed propagation. A flat coherent
pulse kept unchanged in the perturbed metric generally fails this identity.
The same nonminimal action supplies the scalar propagation correction,
the $\xi G_{\mu\nu}\Phi^2$ stress and the Wald coefficient.

This is a conditional leading two-jet stress law, not an exact
finite-radius curved modular identity. An arbitrary bent slice already
gives a false source without its extrinsic correction.
The full quantum reference stress and preparation apparatus must also
enter the same renormalized source. Thus a common action supplies
necessary source and coefficient relations, while the actual curved
modular response remains an additional problem.

\section{The mechanism to be realized}
\label{sec:source-principle}

A physical history must do two jobs at once: form and maintain the
accessible structure, and supply the total source to which the
geometric law responds. The finite return of Section~\ref{sec:physical-return}
establishes a source-to-operation edge inside an assigned gravity
theory. Proposition~\ref{prop:common-source-bridge} establishes what
an all-ball consistency law requires in a different, higher-dimensional
matter theory. Their relation is the following strong target.

\begin{conjecture}[Common-action selection and geometric consistency]
\label{conj:common-action-selection}
There is a local matter--controller theory in which readable histories
form and maintain a nontrivial selected geometric structure, and whose
causal effective action uses the same prepared fields, support and clock
in its source and response. In a low-energy sector with a common cone,
its selected ball data satisfy the geometric consistency condition of
Consciousness~II. Its first-order response is the corresponding Einstein
response with one source normalization; its higher-order response includes
the interaction, contact and clock work that produced the selection.
Selection changes this response through physical couplings and sources,
while a change of description at fixed matter history leaves the source
unchanged. The sustainable family of Conjecture~\ref{conj:gravitating}
is a finite operating sector of this theory, with a response-to-operation
gain that survives the growing-window limit.
\end{conjecture}

The conjecture has tests beyond assigning gravity to a massive apparatus.
Removing the history's access to a command must remove the
record-mediated component of the source, while holding the commanded
interaction fixed and merely discarding a log must preserve it.
Source identification must survive switched physical references;
logical noise records alone fail this test in
Proposition~\ref{prop:weighted-source-ambiguity}.
The same calibrated source must then pass the all-ball consistency
test with fixed initial or boundary data. A law that reproduces only
the JT dilaton or one binary return has not passed the higher-dimensional
test. At second order the source, geometry and changed preparation must
be varied together; entropy Hessians do not become SLD or spacetime
metrics by name.

A model satisfying the conjecture must output a nontrivial selected
sector from a family of retained histories; a prescribed sequence of
couplings is not enough. With a supplied noise and readable seed,
Foundation~I gives a relative stable-hull construction. The next demand
is that the same history dynamics prepares those inputs and that the
resulting region data satisfy Eq.~\eqref{eq:D-MC}, with the total source
and one area coefficient. The scalar surface match removes an
independent matter boundary coefficient, but does not remove MC.
This is the precise gap between common-action source formation and
a selected gravitational law.

The hardest competing explanation is ordinary matter feedback on an
unchanged gravitational stage. It explains every finite JT return here.
To go beyond it, the selected localization family must predict a
common geometric response without refitting the source reference,
area coefficient or region prescription for each history. A permanent
mismatch of that test would refute this common-action realization,
even if its controller remained useful.

\section{From a Selected Interface to a Physical Source}
\label{sec:source}

A calibrated interface acts through a field coupling whose outgoing stress
sources the supplied JT theory. This gives a physical selection--source map;
the origin of the gravitational law in Conjecture~\ref{conj:main} is a
separate question.

\subsection{What records can and cannot do to a source}
\label{sec:records}

Let a classical record $R\sim\mathrm{Bernoulli}(p)$ control whether a physical operation prepares the matter state $\psi$ or leaves the reference $\omega$. Without reading $R$, the matter state is $\rho_p=(1-p)\omega+p\psi$ and its stress is the mixture of the two stresses. Deleting or ignoring $R$ after the preparation does not act on the matter, so $\rho_p$ and its source are unchanged. Changing the preparation probability $p$ changes the source. The information side has a matching statement. For the joint state $\rho_{RM}=(1-p)|0\rangle\langle0|\otimes\omega+p|1\rangle\langle1|\otimes\psi$ and the reference $\pi_R\otimes\omega$, with $\pi_R$ the marginal of $R$, additivity of relative entropy on direct sums and data processing give
\begin{equation}
D(\rho_{RM}\Vert\pi_R\otimes\omega)=p\,D(\psi\Vert\omega),\qquad D(\rho_p\Vert\omega)\le p\,D(\psi\Vert\omega)
\end{equation}
\cite{OhyaPetz1993}. Losing the record lowers the distinguishability available to a reader; it does not remove the matter source. A gravitational consequence of selection must therefore pass through a preparation or interaction, and the conditional geometry of a branch, the average of branch geometries and the geometry of the averaged source are three different objects.

\subsection{The source of a calibrated feedback actuator}
\label{sec:actuator-source}

The companion paper~\cite{Liu2026QRAIFB} realizes a corrective flip of a degenerate bit as a field coupling. A right-moving free massless scalar with null coordinate $u$, current $j=\partial_u\phi$ and stress $T_{uu}=:j^2:$ is coupled through a real smooth switch $g$. We require a finite nonzero one-particle norm
$\lambda^2=\int_0^\infty\omega|\hat g(\omega)|^2d\omega/(4\pi)$ and
$\int(1+u^2)|g'(u)|^2du<\infty$, where
$\hat g(\omega)=\int g(u)e^{-i\omega u}du$ in the $L^2$ sense when needed.
For example, $g\in H^1(\mathbb R)$ together with the weighted derivative condition
is sufficient. Smooth compact switches and the full-line packet below both belong
to this class. The smeared current defines a Weyl unitary,
\begin{equation}
U=\exp\Bigl(-iX\!\int g(u)j(u)\,du\Bigr)=\exp[-i\lambda X(b+b^\dagger)],
\end{equation}
where $b$ is a normalized field mode with mean frequency $\bar\omega$ fixed by $g$, and $\lambda^2\bar\omega=\tfrac14\int g'^2du$. With the drive mode in the coherent state $|\sqrt N\rangle$ and $2\lambda\sqrt N=\pi/2$, the coupling flips the bit with success probability $s_N=(1+e^{-\pi^2/(8N)})/2$. From the current commutator $[j(u),j(v)]=\tfrac i2\delta'(u-v)$ one has $U^\dagger j(u)U=j(u)+\tfrac12Xg'(u)$ and
\begin{equation}
U^\dagger T_{uu}(u)U=T_{uu}(u)+g'(u)Xj(u)+\tfrac14g'(u)^2 .
\label{eq:stress-id}
\end{equation}
The cross term has zero expectation if the bit and the fresh field input are initially independent and $\operatorname{tr}(\rho X)=0$, or if their joint state has no coherence in the bit's $Z$ basis; a $Z$-diagonal bit marginal alone does not suffice~\cite{Liu2026QRAIFB}. Writing $j=\psi(u)b+\psi^*(u)b^\dagger+j_\perp$, the incoming coherent beam has stress $T_{\rm in}=4N[\operatorname{Re}\psi(u)]^2$ and each activation adds
\begin{equation}
F_g(u)=\tfrac14g'(u)^2=4\lambda^2[\operatorname{Im}\psi(u)]^2 .
\label{eq:Fg}
\end{equation}
This added stress is a property of the coupling. Its integral is the switching work $\lambda^2\bar\omega$ supplied by the external switching agent.

The feedback task is that of the Imprint Chain~\cite{Liu2026Chain} with the flip realized by this coupling. A hidden instrument label is estimated from calibration records with mismatch probability $q$; each task step reads the bit with reliability $\nu$, accepts feedback with probability $h$, switches $g$ on when accepted and the reading is negative, and lets a bath with contraction $c$ act. Each opportunity receives its own scheduled incoming beam, independent of the bit and earlier records. In the companion paper the stationary loss $L_\infty(q)$ and the stationary coupling rate satisfy~\cite{Liu2026QRAIFB}
\begin{equation}
p(q)=h\nu L_\infty(q)+\tfrac{h(1-\nu)}2 .
\label{eq:rate}
\end{equation}

\begin{proposition}[Source--loss relation]
\label{prop:source-loss}
At fixed drive, switch and task parameters, the stationary mean source per opportunity outside the apparatus is $\langle T_{uu}\rangle_q=T_{\rm in}+p(q)F_g$. For two calibrations,
\begin{equation}
\langle T_{uu}(u)\rangle_q-\langle T_{uu}(u)\rangle_{q'}=h\nu\,[L_\infty(q)-L_\infty(q')]\,F_g(u).
\label{eq:source-loss}
\end{equation}
\end{proposition}

\begin{proof}
The incoming beam is present at every opportunity. A coupling occurs with stationary probability $p(q)$ and, under the joint initial condition, adds $F_g$ in expectation by Eq.~\eqref{eq:stress-id}. Equation~\eqref{eq:rate} gives the difference.
\end{proof}

Better calibration lowers the task loss and, through it, the source. The coefficient comes from the same interaction that performs the correction; no separate rule maps a calibration error to an emitted pulse. The relation does not assign stress to knowledge. Deleting the calibration record after the commands have been executed leaves the outgoing source unchanged, in line with Section~\ref{sec:records}. The relation is a stationary average per opportunity; a single trajectory, with its fixed hidden match label, does not follow the population mean step by step.

\subsection{A source set by learned quantum protection}
\label{sec:learned-protection-source}

The activation-rate mechanism is not the only way that retained records can
change a source. Brain Anchor~\cite{BrainAnchor} constructs two degenerate
qubits with a programmable coupling
\begin{equation}
 A_a=Z_1+\alpha aZ_2,\qquad
 U_a=e^{-iA_aQ_f},\quad
 Q_f=\int f(u)j(u)du=\sqrt{k}(b+b^\dagger).
 \label{eq:protection-coupling}
\end{equation}
The real smooth compact switch $f$ and positive gain interval are supplied.
The gain $\alpha$ is unknown, and each encounter receives a fresh independent
vacuum field. On $\operatorname{span}\{|01\rangle,|10\rangle\}$ the
coupling is $(1-\alpha a)Z_L$. Calibration records from the commands $\pm h$
have means $r_\pm=e^{-2k(1\mp\alpha h)^2}$, so
$\log(r_+/r_-)=8k\alpha h$. The finite-record certificate in that paper
bounds $|1-\alpha a|$, including command rounding and known drift.
It does not promise an exact fixed algebra after finite training.

\begin{proposition}[One residual controls protection and source]
\label{prop:protection-source}
Fix the acquired command $a$, write $m=1-\alpha a$ and
$F_f=f'^2/4$, and keep the initial field independent of the logical
body and retained ancillas. For every logical input, one encounter has
added mean stress $m^2F_f$ and switching work $m^2E_f$, with
$E_f=\int F_fdu$. After $M$ independent identical encounters, with the
output ports discarded and no interleaved gates or active recovery,
its passive protection error is
\begin{equation}
 \epsilon_M=\tfrac12(1-e^{-2Mk m^2}).
 \label{eq:protection-source-risk}
\end{equation}
For a single pulse with stress moments $E_j=\int u^jF_fdu$ and the
initial data of Proposition~\ref{prop:jt}, the late invariant is
\begin{equation}
 \mathcal C(m)=H_0^2+2\kappa H_0m^2E_1
                 +\kappa^2m^4(E_1^2-E_0E_2).
 \label{eq:protection-source-invariant}
\end{equation}
Thus retained history can change the physical protection and the
gravitational response at fixed pulse shape and activation schedule.
\end{proposition}
\begin{proof}
The current shift is $j\mapsto j+A_af'/2$. The centered independent
field input makes the stress cross term vanish for every logical state,
and $A_a^2=m^2I_L$ on the code. Its characteristic function multiplies
the logical off-diagonal by $e^{-2km^2}$ per encounter; a dephasing
channel's half-diamond error is its phase-flip probability. Time ordering
may give a factor $e^{i\vartheta_f A_a^2}$, which is scalar on the code
but not on the whole two-qubit register. Finally substitute $m^2F_f$
in the fixed-initial-data source solution below.
\end{proof}

This link does not require a rule that emits a pulse only after a wrong
answer. The same interaction produces both the residual error and its
source, even when the logical body is entangled with retained memory.
At $m=0$ the body decouples in the code; a noiseless stored qubit need
not continuously export entropy. The finite-rate entropy-export bound
of Section~\ref{sec:gap} applies to actual renewal, not to every protected
state. Calibration itself adds field energy
$2N(1+\alpha^2h^2)E_f$ for $N$ independent tests at each sign.
Preparations, measurement, the controller and switching mechanism add
further charges. The source suppression during operation alone is not
a total-energy saving.

Brain Anchor also treats an unknown centered Gaussian variance $K_\omega$. Logical records at $0,\pm h$ give
$K_\omega=-\log r(0)/2$ and
$\alpha=[\log r(h)-\log r(-h)]/(8hK_\omega)$; two settings suffice when the signs of $1\mp\alpha h$ are known. Its finite-record mismatch certificate needs no supplied temperature label. The coherence is now $e^{-2K_\omega m^2}$, but the coupling-added source remains $m^2F_f$ on centered independent input.

Incoming stress must still be added. A $u$-clock chiral KMS input has background $\pi/(12\beta^2)$ in natural units. Within a supplied bounded one-temperature family, strictly decreasing $K_\beta$ permits a finite-error background estimate; a general Gaussian variance does not determine incoming energy. At common preparations and initial data, the linear dilaton equation compares coupling-source differences. Equation~\eqref{eq:protection-source-invariant} cannot be used for the full source after dropping its thermal background, nor is $u$ thermality the rapidity-clock reference of Section~\ref{sec:gap}.

The same command can also operate on one reused vacuum field. For nonnegative compact switches $f_j$ on disjoint intervals, with no noncommuting logical gates, the smeared currents commute but have negative correlations. Brain Anchor gives exact coherence $e^{-2m^2k_M}$ with
$k_M=\langle(\sum_jQ_{f_j})^2\rangle\le Mk$, and added source $m^2\sum_jf_j'^2/4$ of energy $Mm^2E_f$. Its JT response uses the actual pulse-train moments and initial coefficients in a valid operating domain. No independent-port approximation, thermal input or infinite-lifetime worldline is inferred.

There is a finite causal pulse in this construction. For example, set
$f(u)=A\exp[-1/(1-(u-3)^2)]$ on $2<u<4$, zero outside, and normalize
its one-particle norm to $k=.2$. Numerical quadrature gives
$(E_0,E_1,E_2)=(.4568522,1.3705565,4.3215400)$.
For $H_0=1$, $\kappa=.05$ and $\alpha=1.1$, the supplied command
$a=48/55$ has residual $m=.04$ and gives
$\mathcal C\simeq1.00021929$; the comparison $a=1$
has $m=-.1$ and $\mathcal C\simeq1.00137053$.
These are deterministic evaluations of the displayed pulse, with code
and outputs in \texttt{supplemental/pulse\_debit}, not statistics of
a calibration pilot. For a random learned command,
the invariant of the mean source uses $(\mathbb E m^2)^2$, whereas
the mean branch invariant uses $\mathbb E m^4$ in the quadratic term. The compact source ends before a later boundary echo is
scheduled, but the calibrated wall and timing instrument remain the
inputs of Section~\ref{sec:echo}.

Stress readback also has an exact blind direction. At fixed $a_0>0$,
the gains $\alpha_\pm=(1\pm d)/a_0$ give residuals $\mp d$ and therefore
the same protection channel and source. Field parity conjugates their
full code interactions while leaving the vacuum and stress invariant.
No parity-covariant stress instrument with invariant outcome labels can
distinguish them, nor can their mean JT responses at common initial data.
Their correcting commands $1/\alpha_\pm$ differ. Probes at $a_0\pm h$
break this ambiguity through the opposite signs of
$8k\alpha h(1-\alpha a_0)$ in their log contrast. A gravitational
response is therefore not a complete calibration reference for the
interface that produced it. Off-frame coupling probes, or a signed
field quadrature, are additional access, with their own sources and costs.
Full output-field parity and a conditional logical pulse can instead
restore the stored state without identifying the gain, as shown in
Ref.~\cite{BrainAnchor}. This recovery leaves the emitted source
unchanged. Active state correction and learning a match that suppresses
the actual source are therefore distinct operations.

The gain is a physical actuator parameter, not a freely changed
description. The selected finite code is not the geometric algebra of
Section~\ref{sec:modular}, and the supplied JT action still supplies the
gravitational response. A link to continuous modular selection requires
that further interface; it is not obtained by identifying a protected
$M_2$ sector with a half-sided modular inclusion.


\subsection{A physical seed changes the source}
\label{sec:projection-source}

The rate-controlled model changes how often a fixed interaction occurs. The coupled
observable can itself change, giving a direct relation between an exactly represented
algebra and a source. Keep the incoming field, switch and activation schedule fixed.
On a qubit put $P_\theta=(I+\cos\theta\,Z+\sin\theta\,X)/2$ and use
$U_\theta=\exp[-i\lambda P_\theta(b+b^\dagger)]$. The projection selects a pointer
algebra and its source population under this fixed schedule.

\begin{proposition}[Selected projection and outgoing stress]\label{prop:projection-source}
For $\lambda\ne0$, an initially independent qubit state $\rho$ and normal field state $\sigma$, the
exactly represented algebra of the induced qubit channel is $\{P_\theta\}'$. Let
$r_\theta=\operatorname{tr}(\rho P_\theta)$. The stress formula below requires its displayed current moments to exist. The outgoing field is
\[
 (1-r_\theta)\sigma+r_\theta e^{-i\lambda Q}\sigma e^{i\lambda Q},\qquad Q=b+b^\dagger,
\]
and its stress is
\begin{equation}
 T_\theta(u)=T_{\rm in}(u)+r_\theta
 \left[g'(u)\langle j(u)\rangle_\sigma+\tfrac14g'(u)^2\right].
\label{eq:projection-source}
\end{equation}
Thus two selected projections can change the source at fixed global input state and
fixed activation schedule.
\end{proposition}
\begin{proof}
The quadrature has non-atomic spectral distribution in a normal state, so the
multiplicative-domain theorem of Q-RAIF~C gives $\{P_\theta\}'$~\cite{QRC}.
Expanding $U_\theta$ in its two projection blocks and tracing the qubit gives the
field mixture; off-diagonal blocks have zero trace even when $\rho$ does not commute
with $P_\theta$. In the current identity replace the Pauli $X$ by $P_\theta$ and
use $P_\theta^2=P_\theta$. Initial independence gives~\eqref{eq:projection-source}.
\end{proof}

For $\rho=(I+0.8Z)/2$, the seeds $P_0$ and $P_{\pi/2}$ select different pointer
algebras, with $r_0=0.9$ and $r_{\pi/2}=0.5$. In the worked coherent packet,
the cross profile is $8\lambda\sqrt N\operatorname{Re}\psi\operatorname{Im}\psi$.
Its first three moments vanish, so the invariant increment is~\eqref{eq:dC} with
$p$ replaced by $r_\theta$. With the table's $N=16$ and background it is
$2.0807581\times10^{-5}$ versus $1.1559770\times10^{-5}$. The full source also
contains the cross profile, which enters the local dilaton even though its first
three moments do not enter the final invariant.

Here the source weight is derived from the state's population in the selected
projection; it is not assigned to an external mismatch label. Changing the projection
means changing a physical interaction, not renaming an unchanged algebra. The two
algebras are abstractly isomorphic but are differently embedded relative to the fixed
state and coupling. This supplies a seed--interface--source link. It does not identify
that finite-dimensional pointer algebra with a Type~III geometric candidate or explain
why the coupling selects this seed rather than another. Those are the remaining
selection and geometry bridges.


\subsection{Response of a given Jackiw--Teitelboim model}
\label{sec:jt}

Take the JT action~\cite{AlmheiriPolchinski2015,Moitra2020}
\begin{equation}
I=\frac1{16\pi G}\int\sqrt{-g}\,\Phi(R+2/L^2)+I_{\rm matter},\qquad
ds^2=-\frac{4L^2\,du\,dv}{(1+uv)^2},
\end{equation}
with the usual boundary terms, $\kappa=8\pi G$, $D=1+uv$ and the vacuum dilaton $\bar\Phi=H_0(1-uv)/D$, $H_0>0$. The field equations are $\nabla_\mu\nabla_\nu\Phi-g_{\mu\nu}\Box\Phi+g_{\mu\nu}\Phi/L^2=-\kappa T_{\mu\nu}$, and in these coordinates
\begin{equation}
\Phi_{uu}+\tfrac{2v}D\Phi_u=-\kappa T_{uu},\qquad
\Phi_{vv}+\tfrac{2u}D\Phi_v=-\kappa T_{vv},\qquad
\Phi_{uv}+\tfrac{2}{D^2}\Phi=\kappa T_{uv} .
\label{eq:jt}
\end{equation}
The matter is a conformal theory with matched left and right central charges, and the state-independent anomaly is removed by subtracting the same background. The left-moving input is unchanged, and the state-dependent source is right-moving, $T_{uu}=F(u)$ with $T_{uv}=T_{vv}=0$.

\begin{proposition}[Response with fixed initial data]
\label{prop:jt}
Let $F$ be integrable with $\int|u|^2|F|du<\infty$. On a connected coordinate domain with $D\ne0$, write the perturbation as $\Phi-\bar\Phi=(A(u)+vB(u))/D$ and impose the initial coefficient data $A(u),B(u)\to0$ as $u\to-\infty$. For a source bounded below in support this is equivalent to zero perturbation before the source. The unique solution with these data is
\begin{equation}
\Phi=\bar\Phi+\frac{A(u)+vB(u)}D,\qquad
A=-\kappa\int_{-\infty}^{u}(u-z)F(z)dz,\qquad
B=-\kappa\int_{-\infty}^{u}z(u-z)F(z)dz .
\label{eq:retarded}
\end{equation}
The perturbation is linear in $F$. With $E_j=\int u^jF\,du$, after the source (or in the asymptotic limit for sources with tails) the invariant $\mathcal C=\Phi^2-L^2(\nabla\Phi)^2$ takes the vacuum value
\begin{equation}
\mathcal C=H_0^2+2\kappa H_0E_1+\kappa^2(E_1^2-E_0E_2),
\label{eq:invariant}
\end{equation}
and inside the support $\partial_u\mathcal C=\kappa F\,(2H_0u+uA-B)$.
\end{proposition}

\begin{proof}
Write $\psi=D\,(\Phi-\bar\Phi)$. The $vv$ equation of~\eqref{eq:jt} reads $\psi_{vv}=0$, so $\psi=A(u)+vB(u)$, and the other two equations reduce to $A''=-\kappa F$, $B''=uA''$ and $B'=uA'-A$, which Eq.~\eqref{eq:retarded} satisfies. The absolute first and second moments make both retarded coefficients tend to zero at $-\infty$. Two solutions differ by $\alpha+\beta u+\gamma v-\alpha uv$, whose coefficients are $A_h=\alpha+\beta u$ and $B_h=\gamma-\alpha u$. Their two zero limits force $\alpha=\beta=\gamma=0$. Every vacuum solution has numerator $\alpha+\beta u+\gamma v-\alpha uv$ and invariant $\alpha^2+\beta\gamma$; after the source $\alpha=H_0+\kappa E_1$, $\beta=-\kappa E_0$, $\gamma=\kappa E_2$. The balance law follows by differentiation. 
\end{proof}

The coefficient condition cannot be replaced by $\Phi-\bar\Phi\to0$ at each fixed $v$. The homogeneous solution $\eta v/(1+uv)$ has that limit for every fixed $v$ along nonsingular points. It even tends to zero uniformly on the subdomain $u<0$, $v\le0$, since its absolute value is at most $|\eta|/|u|$; no uniform limit near $D=0$ is implied. Adding it leaves the source unchanged but changes the late invariant by $-\kappa E_0\eta$. It is excluded by $B(-\infty)=0$. Source comparisons below hold $H_0$, the coordinate matching and both coefficient boundary data fixed.

The JT curvature stays $-2/L^2$; the dilaton and the invariant respond. A dilaton difference is not by itself a higher-dimensional curvature or a mass, and comparing boundary proper times needs a boundary normalization.

The matching condition matters. Set $A_F,B_F\to0$ at $+\infty$, or exactly to zero after a compact source, instead of the initial data of Proposition~\ref{prop:jt}. The final-state-matched solution differs from the retarded solution by
\[
A_F=A_R+\kappa(uE_0-E_1),\qquad
B_F=B_R+\kappa(uE_1-E_2).
\]
For the chiral diffeomorphism states of Consciousness~I~\cite{PaperCI}, with a Cut at $a$ to the left of the support, the relative entropy is $D_a=2\pi\int(u-a)F(u)\,du$. The final-state-matched solution gives $\delta\Phi(a,0)/(4G)=-D_a$, while the initial-state-matched solution gives $\delta\Phi(a,0)=0$ at a Cut before the pulse. The first identity compares an earlier section with a matched final state; it does not say that a pulse removes area before it arrives. The dilaton--modular-energy relation itself is known~\cite{Callebaut2019}; what is used here is its type-correct pairing with a specified source and matching condition.

The same equations show that a null-only record of the source is incomplete. A genuine vacuum-density shift $T_{\mu\nu}=-\delta\rho\,g_{\mu\nu}$ has zero null components but gives the particular solution $\delta\Phi=\kappa L^2\delta\rho$. A probe that reads only null energy cannot certify that a source is absent, which is the source-completeness condition relevant to the vacuum curvature test of Ref.~\cite{QRNV}.


\section{Response that changes the next operation}
\label{sec:physical-return}

A source-to-geometry calculation becomes a return only when the geometry
changes an actual record and that record changes a later interaction.
The readout permission matters. In the minimally coupled JT sector,
instruments on the same matter metric with the same inputs have identical
adaptive record laws even when the dilaton differs
(Proposition~\ref{prop:blind-readback}). A boundary tied to a fixed dilaton
value evades that obstruction. Its echoes reconstruct the three vacuum
coefficients, and a source-compensated inverse includes finite probe
backreaction (Propositions~\ref{prop:echo} and
\ref{prop:compensated-echo}). An exact current contact instead obtains an
invariant-dependent bit, including the contact's own source
(Propositions~\ref{prop:driven-boundary-sensor} and
\ref{prop:weak-boundary-sensor}). The bit selects the next physical
projection. Its same-field conditional covariance remains part of the
history; detector replacement does not reset the field.

These are finite physical records, not permission to read an exact metric.
The worked current contact has total variation only
$2.3650\times10^{-4}$ between its two initial geometries.
Its fixed-cutoff energy budget also forbids indefinitely many contacts of
unchanged positive source charge. Appendix~\ref{app:boundary-return}
contains the timing, stress and accumulated-source bounds.

\subsection{An acquired record drives a new material source}
\label{sec:record-actuated-source}

A record can act without first being converted into a supplied mechanical
initial condition. Keep the two classical scalar fields and the material
action of Section~\ref{sec:inertial-scalar-support}, and append
$i\langle\Psi,\dot\Psi\rangle-\langle\Psi,H_c\Psi\rangle$ to its
proper-time Lagrangian. The normalized state $\Psi$ contains a finite
history clock, probe $P$, record $R$, replacement blank $B$ and an old
purification $Q$. Set
\begin{equation}\label{eq:record-spring-core}
 H_c(\Phi,y)=H_{\rm hist}+E_{\rm off}\mathbf1
       +\lambda(\Phi-\phi_0)Z_P
       +\frac{\chi_R}{2}P_R(y-d)^2,
 \qquad P_R=|1\rangle\langle1|_R,\quad\chi_R\ge0.
\end{equation}
$H_{\rm hist}$ is the positive gate-conjugated clock of Q-RAIF~C.
Its first nonidentity edge performs the $Y_P$ acquisition $C_Y$, the next performs a true
SWAP of $P$ and $B$, and the remaining edges are identities.
There is no proper-time switch. Initially $R=0$, so the spring has
zero expectation. The actual acquisition opens its force channel.
$H_c\ge0$ while $E_{\rm off}>|\lambda(\Phi-\phi_0)|$.
The target $d$, couplings and initial clock are prepared hardware.
The scalar fields, mechanical coordinates and geometry use an Ehrenfest
mean-action approximation~\cite[Section~2]{CastroGross2013};
the finite quantum core and its old reference are retained
in full, without sampling or collapsing $R$.

Here $\theta$ is a material angle, $y$ its storage coordinate,
$I,J,k,K>0$ their inertia and elastic parameters, and $U(\Phi)$
the material support potential. Proper-time derivatives obey
$I\ddot\theta+k(\theta-y)=-\gamma\dot\theta$,
$\gamma=g_+^2+g_-^2$, with no incoming classical waves.
Their emitted powers are $P_\pm=g_\pm^2\dot\theta^2$.
The mechanical stored energy is
$e=I\dot\theta^2/2+J\dot y^2/2+k(\theta-y)^2/2+Ky^2/2$,
and $w=\dot\Phi$ along the wall.
The positive wall junction and its two invariant charges are specified
by Eqs.~\eqref{eq:material-junction} and \eqref{eq:two-outgoing-wall};
all of these mechanical and gravitational data are supplied.

Let $p=\langle P_R\rangle$, $z=\langle Z_P\rangle$ and
$E_c=\langle H_c\rangle$. Variation of this one action gives
\begin{align}\label{eq:record-spring-work}
 i\dot\Psi&=H_c\Psi,& F_y&=\chi_R(y-d)p,& F_\Phi&=\lambda z,\\
 J\ddot y+k(y-\theta)+Ky&=-F_y,&
 \dot e&=-P_+-P_--F_y\dot y,&
 \dot E_c&=F_y\dot y+F_\Phi w.\nonumber
\end{align}
The angle equation and $P_\pm=g_\pm^2\dot\theta^2$ are unchanged.
Here $p$ is an expectation obtained by variation, not a decoded number
fed to a separate control policy. The mass and its tangent balance are
\begin{equation}\label{eq:record-spring-mass}
 M=m_{\rm rest}+U+e+E_c,\qquad
 \dot M=(U'+F_\Phi)w-P_+-P_- .
\end{equation}
Consequently Eq.~\eqref{eq:two-outgoing-wall} holds with $U'$ replaced
by $U'+F_\Phi$ in its acceleration equation; both invariant and null-clock
equations retain their directions. Differentiating the mass junction
with these equations gives zero. The separate work balances are
\begin{align}\label{eq:record-spring-ledgers}
 e(\tau)+Q_{\rm out}(\tau)+\int_0^\tau F_y\dot y\,ds&=e(0),\\
 E_c(\tau)-E_c(0)&=\int_0^\tau(F_y\dot y+F_\Phi w)\,ds.\nonumber
\end{align}
The clock interaction and conditional spring are included in $E_c$.
The spring breaks the symmetry of Proposition~\ref{prop:force-neutral-meter}; $F_\Phi$ cannot
be set to zero in this coupled calculation.

\begin{proposition}[A nonzero record-mediated source return]
\label{prop:record-mediated-source}
Take two initial geometries with the same $\Phi=\phi_0$, $w=0$, mass,
mechanical data and quantum input, but different $\mathcal C_-$;
set $\mathcal C_+$ by the mass junction. In the positive exterior domain,
either $\lambda=0$ or $\chi_R=0$ makes their proper emitted powers
identical. With both interactions present, there are positive finite-clock
data for which the emitted powers differ on a nonempty finite interval.
Their record, mechanical response, source and geometry obey the same action
and work account.
\end{proposition}
\begin{proof}
For $\chi_R=0$, the mechanical subsystem is independent of the reader
and geometry. For $\lambda=0$, the joint mechanical and quantum equations
are independent of geometry. Either condition fixes identical mechanical
sources for the two initial geometries.

For a nonzero example choose $\theta(0)=\dot\theta(0)=0$,
$y_0(y_0-d)\ne0$ and probe state $(I+\upsilon X_P)/2$, $\upsilon\ne0$,
with a blank $R$ and its old purification held. Let $r\ge1$ be the
$C_Y$ edge, $a_j>0$ the clock edge strengths and
$G_r=\prod_{j=1}^r a_j$. The initial acceleration difference
$\delta a$ follows from the junction equations, so that
$\delta\Phi(t)=\delta a\,t^2/2+O(t^3)$.
The shortest-path Dyson expansion gives
\begin{equation}\label{eq:record-spring-leading-record}
 \delta p(t)=K_r t^{2r+3}+O(t^{2r+4}),\qquad
 K_r=\frac{\upsilon G_r^2\lambda\delta a}
             {6r!(r-1)!(r+3)}.
\end{equation}
An insertion of $\lambda\delta a\,t^2Z_P/2$ after $j<r$ forward
edges contributes the time integral
$t^{r+3}(j+1)(j+2)/(r+3)!$.
Their sum is $t^{r+3}r/[3(r+3)r!]$.
Using $C_Y^*P_RC_Y=(I+Y_P)/2$ in the interference with the unperturbed
shortest path gives the coefficient above. The insertion after all
$r$ edges has Hermitian $P_RZ_P$ in that interference and zero real
part. Further graph paths and spring insertions raise the order;
the spring annihilates the blank record before acquisition.

Integrating the mechanical equations then yields
\begin{equation}\label{eq:record-spring-leading-source}
 \delta(P_++P_-)(t)=
 -\frac{2\gamma k^2 y_0\chi_R(y_0-d)K_r}
 {I^2J(2r+4)(2r+5)(2r+6)}t^{2r+7}+O(t^{2r+8}).
\end{equation}
The baseline angle velocity is $ky_0t/I+O(t^2)$.
Different initial $\mathcal C_-$ give a nonzero $\delta a$ for generic
supported data, since initially
$\dot w=-U'\beta_+\beta_-/M-\phi_0/L^2$.
The displayed coefficient is nonzero for the stated nonzero parameters.
The coupled finite-dimensional equations are analytic on their strict
positive exterior domain. Local existence and continuity therefore give
a nonempty admissible family, not just a numerically resolved difference.
\end{proof}

The proof isolates the causal mechanism analytically, rather than
attributing a small numerical difference to geometry. Blocking either
sensing or record actuation removes the source contrast. Keeping the
record but blocking actuation preserves a readout difference while
removing the later source difference. Its numerical realization and
complete mass/work checks are in Appendix~\ref{app:material-return}.
The two histories here differ in initial gravitational data; they do
not purport to generate those data from algebra selection.

For histories with common initial geometric constraints, the distinct
spatial construction of Proposition~\ref{prop:finite-spatial-return}
has a nonzero source--dilaton--sensor path and a later material response.
Its finite-width regulator, support and synchronization are physically
varied. The complete material response also contains ordinary clock
and force effects; only the separately isolated path is gravitational.
A third construction, Proposition~\ref{prop:covariant-quartet-bulk},
couples the actual full-mode quartet state to paid clock/director
fields and constrained JT bulk evolution. It supplies a local Ward
identity and a causal prepared bulk domain, not moving-wall gate
matching or the full thermal apparatus.

A finite record error alone does not control an unbounded stress.
The response comparison uses weighted source bounds
(Proposition~\ref{prop:finite-prefix-response}), and the coupled
trajectory--reader bound requires domain-preserving strong solutions
and the stated bounded-coupling derivative assumptions
(Proposition~\ref{prop:coupled-return-tolerance}).
Appendix~\ref{app:finite-return-details} gives these estimates and the
separate full-mode clock-work matching result. None converts
fixed-window existence into sustained operation.

\section{Maintenance and its gravitational price}
\label{sec:gap}

\subsection{Entropy removal and the source it can constrain}
\label{sec:maintenance}

A fixed unitary correction can return its control port coherently. Removing entropy
injected by unobserved noise is a different task. To price it, the state whose entropy
is removed must include the controller's usable side information. The heat account
must also distinguish a thermal reservoir from a supply of pure memories.

HAFF~G supplies a record-driven program and executor in which selecting
a physical coupling changes its source under one fixed executor
\cite{Liu2026PaperG}. The dephased quartet below adds a specified full
renewal debit. It does not identify the two material realizations or
inherit a thermal controller from that program loop.

\begin{proposition}[Maintenance entropy account]\label{prop:entropyaccount}
Let $V=SC$ include the maintained system and every retained controller and
side-information register. A correction acts unitarily on finite-dimensional
$VAB$, initially in $\rho_V\otimes\alpha_A\otimes\tau_B$, where
$\tau_B=e^{-\beta H_B}/Z$, $\beta=(k_BT)^{-1}>0$. Put
\[
 s=S(V)-S(V'),\qquad p=S(A')-S(A),\qquad
 Q=\operatorname{tr}H_B(\rho'_B-\tau_B).
\]
Then, with $I_{V:A:B}=S(V)+S(A)+S(B)-S(VAB)$,
\begin{equation}
 \beta Q+p=s+I_{V':A':B'}+D(\rho'_B\Vert\tau_B)\ \ge s.
 \label{eq:entropyaccount}
\end{equation}
If the marginal entropy of $C$ is returned, $s$ equals the decrease of
$S(S|C)$. In a sequence of such corrections, a bounded total purity consumption
and a positive asymptotic entropy-removal rate $\dot s$ imply
$\liminf Q_{[0,t]}/t\ge k_BT\dot s$.
\end{proposition}
\begin{proof}
Unitary invariance of total entropy and the initial product give
$\Delta S(V)+\Delta S(A)+\Delta S(B)=I_{V':A':B'}$.
Also $\beta Q=\Delta S(B)+D(\rho'_B\Vert\tau_B)$.
Combine the identities and sum over corrections. The rate statement follows when
the accumulated $p$ divided by elapsed time tends to zero.
\end{proof}

This is the Landauer entropy balance, with the memory-purity term displayed
explicitly; the equality and the role of accessible side information are standard
\cite[Theorem~3 and Section~5.1]{ReebWolf2014}. For initially correlated or
nonthermal inputs the same calculation instead gives
\[
 \beta Q+p=s+(I_{\rm out}-I_{\rm in})+
 D(\rho'_B\Vert\tau_B)-D(\rho_B\Vert\tau_B).
\]
These consumed correlations and bath free energy can pay for correction. A classical
copy of the error permits a reversible conditional flip with $s=0$, although the
system's marginal entropy decreases. Negative quantum conditional entropy gives
another reason to include the side information. A pure-memory swap has $p=s$ and
can have $Q=0$. A retained oscillator's closed unitary loop also has no entropy
removal by itself. Thus a count of corrections or a marginal noise entropy is not
automatically $\dot s$. In a cyclic device $\dot s$ must be established from its
actual noise and controller dynamics. For nondegenerate systems, heat and supplied
work are further separated by the change in internal energy.

\begin{proposition}[A finite-rate chiral exhaust floor]\label{prop:chiral-floor}
Consider a unitary, zero-chemical-potential chiral conformal exhaust with central charge
$c_{\rm ex}>0$, thermal input at $T\ge0$, and an extensive entropy-current account.
Energy and temperature are defined by the same translation generator and clock.
Assume that the outgoing modes export net dimensionless entropy current
$x\ge\dot s>0$, where $\dot s$ is the actual maintenance entropy-removal rate after
all retained side information and sublinear resource consumption are accounted for.
The entropy-current maximum at fixed energy is the thermal one. Then the excess
heat current over the incoming thermal stream obeys
\begin{equation}\label{eq:chiral-floor}
 \dot Q\ge k_BT\dot s+\frac{3\hbar}{\pi c_{\rm ex}}\dot s^2.
\end{equation}
\end{proposition}
\begin{proof}
Write $a=\pi c_{\rm ex}k_B^2/(12\hbar)$. The thermal energy and dimensionless
entropy currents are $J_E(T)=aT^2$ and $J_S(T)=2aT/k_B$.
The first is the standard chiral CFT current~\cite[Eq.~(1)]{BernardDoyon2012};
the second follows from its extensive Gibbs entropy. Entropy maximization at
fixed energy is the Gibbs variational principle, underlying the single-channel
information-flow bound~\cite{Pendry1983}.
Choose $T_*=T+k_Bx/(2a)$ so that the thermal output has the required entropy.
Any output with at least that entropy has energy current at least $aT_*^2$.
Subtracting $aT^2$ gives
\[
 \dot Q\ge k_BTx+\frac{k_B^2x^2}{4a}
            =k_BTx+\frac{3\hbar x^2}{\pi c_{\rm ex}}.
\]
The expression increases for $x\ge0$.
\end{proof}

The current account is a physical condition on the exported modes, not the entropy
of a Type~III local interval. Equation~\eqref{eq:joint-exhaust} establishes the
needed joint entropy rate in a fresh-port account; matching those ports to a conformal current
must preserve its capacity and entropy budget. At fixed capacity and positive
rate the quadratic term survives even at $T=0$. Thermal output saturates the
channel inequality, but this alone does not construct a correcting apparatus.
The currents are long-window, extensive Gibbs limits. For an observation window
$\tau$ with net exported entropy $X_\tau$, bounded rates and sublinear boundary
corrections give
\[
 Q_\tau\ge k_BT X_\tau+
 \frac{3\hbar X_\tau^2}{\pi c_{\rm ex}\tau}+o(\tau).
\]
The current formula is not an exact inequality for an arbitrarily short finite
packet. Applying the Gibbs bound to the entire output also includes correlations
between packets; it does not require each output to be a separate thermal state.

\subsection{A conditional gravitational reading}

The price becomes a geometric reading only with a physical transport map.
The following JT coefficient reads the total routed heat in one
matched clock and reference. It does not assign curvature to entropy
that remains in a memory or exits through another channel.

\begin{proposition}[A heat account read through the dilaton]\label{prop:heatslope}
Compare two JT solutions with the same $H_0$, boundary coefficient data and
coordinate and clock normalization. Suppose their source difference $F$ is the
exhaust of the correction in Proposition~\ref{prop:entropyaccount}, with finite
absolute second moment and $\int Fdu=Q$. This last equality requires that the
specified outgoing channel carry all the heat in the same energy normalization.
Then its retarded coefficient difference obeys
\begin{equation}
 -\frac{A'(+\infty)}\kappa=Q\ \ge\ k_BT(s-p).
 \label{eq:heatslope}
\end{equation}
For a compact exhaust and any $u_2>u_1$ after it, the left side is exactly
$-[A(u_2)-A(u_1)]/[\kappa(u_2-u_1)]$. No pointwise sign of $F$ is required.
\end{proposition}
\begin{proof}
Differentiate Eq.~\eqref{eq:retarded}: $A'(u)=-\kappa\int_{-\infty}^uF(z)dz$.
Take the late limit and apply Eq.~\eqref{eq:entropyaccount}.
\end{proof}

The reference for $Q$ is the incoming thermal reservoir transported through the
same exhaust, before its heat uptake; $H_B$ is its $u$-translation generator.
A thermal background has a nonintegrable stationary flux, so the JT solution here
is the difference from that same transported thermal background.
The identification $E_0=Q$ is an energy-transport hypothesis:
other exhaust channels, energy stored in the device, work beams crossing the
observation surface and redshift must be accounted for separately. With a fixed
transport factor $E_0=\zeta Q$, Eq.~\eqref{eq:heatslope} acquires the same factor
$\zeta$. The coefficient slope is read from the dilaton in the declared matched
coordinate frame, not an invariant under arbitrary changes of boundary clock.
The relation links an actual entropy-removal process to a specified dilaton reading.
The JT dynamics and heat-routing map are its geometric inputs.


For a matched running single-chiral source, the marginal-horizon
balance is $dS_h/d\eta=2\pi F\xi^2$. If the charged heat is exactly
$F\xi^2$ and the physical input is thermal at $T_H=1/(2\pi)$ in the
same rapidity clock, the long-window bound becomes
$\overline{\dot S_h}\ge\dot s+6\dot s^2/c_{\rm ex}$.
The actual thermal-source and reference-clock terms, positive-domain
conditions and nonthermal contact excess are retained in
Appendices~\ref{app:maintenance-details} and \ref{app:running-clocks}.
A reparametrized vacuum is not a prepared heat reservoir.

\section{Selection and response over growing windows}

\subsection{A conjecture on a sustainable selection--response loop}

The unresolved question is whether one reused local body can combine
selection, renewal, sharp exhaust and informative geometric feedback.
Fix two physical preparations on a logical receiver $A$ and its current
purifier $R$:
\[
 |\Phi_Z\rangle=\sqrt{1-\epsilon}|00\rangle+
                   \sqrt\epsilon|11\rangle,\qquad
 |\Phi_X\rangle=(H_A\otimes I_R)|\Phi_Z\rangle,
 \quad 0<\epsilon<1/2.
\]
The logical labels denote Q-RAIF~C's neutral pair codes, with their
spectator sea retained. Here $R$ is the paid, resettable current purifier
in the third and fourth body species, available to joint neutral gates.
Old references outside this source quartet are denoted $R_{\rm old}$
and are not acted upon. The quartet is a resettable working source
alongside the series' protected unknown content $H$, not a replacement
for it; access to the current $R$ comes from those paid body species
and joint neutral gates, not from an old-purifier permission. Each cycle's inaccessible environment applies
the same pointer-aligned channel
\[
 \mathcal D_q(\rho)=(1-q)\rho+qZ_A\rho Z_A,
 \qquad 0<q<1/2,\qquad h(x)=-x\ln x-(1-x)\ln(1-x).
\]
No noise error label is available to the apparatus. Renewal returns the
complete quartet, including the current purifier.

\begin{proposition}[Pointer alignment and the complete renewal debit]
\label{prop:pointer-debit}
For $\rho_p=\mathcal D_q(|\Phi_p\rangle\langle\Phi_p|)$, return to its
specified pure quartet removes entropy $h_p=S(\rho_p)$, where
\begin{equation}\label{eq:pointer-debit}
 h_Z=h\left(\frac{1-\sqrt{1-16q(1-q)\epsilon(1-\epsilon)}}2\right)>0,
 \qquad h_X=h(q)>h_Z.
\end{equation}
The receiver-only debits are $0$ for $Z$ and
$h(\epsilon+q(1-2\epsilon))-h(\epsilon)$ for $X$.
Approximate quartet return subtracts its residual joint entropy from
$h_p$, with finite-dimensional entropy continuity controlling the error.
Program updates, sensors and usable controller registers enter the full
maintenance debit in Proposition~\ref{prop:entropyaccount} as well.
\end{proposition}
\begin{proof}
The mixture of $|\Phi_p\rangle$ and $Z_A|\Phi_p\rangle$ has nonzero
eigenvalues $(1\pm d_p)/2$, with
$d_p^2=1-4q(1-q)(1-|\langle\Phi_p|Z_A|\Phi_p\rangle|^2)$.
The overlaps are $1-2\epsilon$ and zero. Strict monotonicity of $h$ on
$(0,1/2)$ proves the contrast. Tracing $R$ gives the receiver debits.
The reference marginal is unchanged by the noise, so the difference
between each complete and marginal debit equals the receiver--purifier
mutual information lost to the noise.
\end{proof}

Pointer alignment lowers the complete renewal price without making it
zero. For $\epsilon=.1,q=.05$,
$h_Z=.0877516288$ and $h_X=.1985152433$ nats per cycle.
The marginal $Z$ state is unchanged, but its purifier correlations are not.
The loop must establish that its additional charged operations preserve
an actual complete-debit contrast $\bar s_X>\bar s_Z$.

\subsection{Three distinct targets}
\label{sec:sustainable}

The finite return is target A: an actual acquisition changes a source,
its response changes another acquisition, and the retained history
remains in the same finite mean-action evolution.
Target B is sustained thermal renewal of the body, controls and ports
over growing runs while the independent unknown content stays held.
Target C combines that renewal with a sharp complete exhaust and
nonvanishing causal response gain. A does not prove B, and B alone
does not prove C.

Conjecture~\ref{conj:gravitating}, stated in full in
Appendix~\ref{app:sustainable-details}, makes C quantitative.
At fixed cycle rate and exhaust capacity, it requires complete-output
sharpness, normalized whole-window response contrast, and a positive
average gain at \emph{all} declared decision opportunities on the
actual-run history measure. Every sensor, program update, delay,
buffer, work bank and current purifier enters the account; old
inaccessible references are not reset. Resource size may grow with
the run, while uncharged purity and correlation consumption have no
positive asymptotic rate.

This target has real obstructions. Accurate direct purifier dumping
retains cross-species redundancy and cannot make the complete Gibbs
deficit vanish. Individually sharp program phases still pay a positive
Jensen variance cost when their entropy rates differ at nonzero duty
fractions. An exploration scheme whose only response updates occur
on a vanishing fraction of opportunities has vanishing average gain,
even with infinitely many such updates. These facts do not exclude
every joint protocol; they exclude tempting combinations of successful
components. The complete statements and resource clocks remain in
Appendix~\ref{app:sustainable-details}.

The following fixed-source converter supports one part of the target.
It does not supply the adaptive source's product structure.

\begin{proposition}[Finite-carrier sorting of the dephased source]
\label{prop:pointer-sort}
Q-RAIF~C's charge-sector sorter~\cite{QRC} applies to $L$ independently
realized copies of either $\rho_p$ under its finite-carrier hypotheses:
the source marginal is $\rho_p^{\otimes L}$, initially a product with
the complete thermal carrier; their external purifications are retained;
joint source operations act on both the receiver and its current $R$
pair, while identity on old references refers to systems outside that
quartet;
the source is neutral in the original zero-mode quartet sector; the
original neutral gates and paid loading, unloading and pipeline stages
are available; and weighted energy tails control the implementation
limit. It returns the source near its blank with entropy slack $o(L)$
and complete output Gibbs deficit $o(V)$, $V=L/\nu$, at
\[
 T_{*,p}=T_H+\frac{6\nu h_p}{\pi c_{\rm ex}},\qquad c_{\rm ex}=4.
\]
The same spectral argument applies to a supplied deterministic word
with product marginal $\bigotimes_{j=1}^L\rho_{p_j}$ and actual entropy
$s_L=\sum_jh_{p_j}$. The finite-bank target uses $s_L$ plus the same
subextensive margin. If $s_L/L\to s$, its temperature tends to
$T_H+6\nu s/(\pi c_{\rm ex})$, with no convergence rate required;
otherwise there need not be one limiting temperature. This is a block
converter, not an adaptive moving apparatus or a per-cycle heat-reading law.
\end{proposition}
The spectral proof is in Appendix~\ref{app:source-clock-support}.
It uses a supplied product marginal and controlled energy tails, not
a local reset certificate for an adaptive history.

\section{Relation to Existing Approaches}
\label{sec:relation}

AdS/CFT~\cite{Maldacena1999} and its area--entropy relation~\cite{RyuTakayanagi2006} supply the dictionary used in Theorem~\ref{thm:linearized}; it is not derived here or granted to arbitrary selected algebras. MERA~\cite{Vidal2008,Swingle2012} implements coarse-graining through disentanglers and isometries, with network depth resembling a radial AdS direction. Different networks can describe the same boundary data. Relating either construction to physical selection requires the explicit algebra-to-geometry map.

Jacobson's horizon Clausius argument~\cite{Jacobson1995} shares the entropy-based approach to gravitational dynamics. Section~\ref{sec:source} addresses its source-side input: which selecting processes produce the stress whose response is computed. Loop quantum gravity~\cite{Rovelli2004,Thiemann2007} instead quantizes the metric with background independence. Its relation to a path of selected algebras would require a representation map from spin-network states. These connections motivate the proposal without supplying its missing selection law.

\section{Conclusion}
\label{sec:conclusion}

The gravitational proposal is that physical histories form a localization
structure whose slow evolution admits a common geometric law.
Modular stability does not select that structure uniquely.
A specified half-sided modular family supports conditional holographic
matching; the conformal scalar bridge instead relates one prepared
source and its Wald surface coefficient to an all-ball test in the
same spacetime. Nonzero-frequency preparations can identify the
intrinsic response within its stated class. The consistency condition
itself still needs a physical explanation.

The source mechanism is explicit. Changed correction, protection and
projection interactions change stress. A finite JT mean-action model
retains the acquired quantum record, old reference, mechanics, clock,
support and recoil, and proves a nonzero record-mediated source return.
The spatial return and full-mode quartet bulk strengthen different
interfaces. They are evidence for a common theory, not one already
joined higher-dimensional apparatus.

Maintenance can favor one selected structure over another.
Pointer alignment lowers the complete renewal debit, and a thermal
chiral channel of fixed capacity places a quadratic rate price on
its export. Matched horizon transport can read that price, provided
the actual source, clock reference and heat routing are supplied.
Sharp conversion of a fixed product source does not establish
sustained renewal or sharp informative adaptive return.

The remaining mechanism is therefore specific: histories must form
the accessible sector, the selecting interactions must carry their
reaction stress, and one calibrated source must satisfy the same
geometric consistency test across those histories. At second order
the geometry, propagation, source and region prescription must vary
together. A visibility filter cannot replace this dynamics, and
changing an algebra at fixed stress cannot create a second
first-order sourced Einstein tensor. Establishing this common
selection--source--geometry law would turn the adiabatic algebra
proposal from an interpretation of supplied gravity into a mechanism
for its effective dynamics.

